\section{Formulação do VRPTW e Heurística Solomon I1}
\label{sec:vrptw_i1}

\subsection{Modelo do problema}

Considere um grafo completo $G=(V,E)$ com $V=\{0\}\cup N$, onde $0$ é o depósito e
$N=\{1,\ldots,n\}$ é o conjunto de clientes.
Cada cliente $i\in N$ possui:
\begin{itemize}
  \item demanda $q_i$;
  \item janela de tempo $[a_i,b_i]$;
  \item tempo de serviço $s_i$.
\end{itemize}

A frota é homogênea, com capacidade $Q$ e no máximo $m$ veículos.
No padrão Solomon, a distância entre dois nós $i$ e $j$ é dada por distância euclidiana truncada em uma casa decimal \cite{solomon1987algorithms}:
\[
 d_{ij}=\left\lfloor 10\cdot \sqrt{(x_i-x_j)^2+(y_i-y_j)^2}\right\rfloor/10.
\]

Para uma rota $r=(0,r_1,\ldots,r_k,0)$, o cronograma é calculado de forma recursiva:
\begin{align}
 t^{arr}_{r_h} &= t^{dep}_{r_{h-1}} + d_{r_{h-1},r_h}, \\
 t^{start}_{r_h} &= \max\{t^{arr}_{r_h}, a_{r_h}\}, \\
 t^{dep}_{r_h} &= t^{start}_{r_h} + s_{r_h}.
\end{align}

Uma solução é factível se:
\begin{enumerate}
  \item cada cliente é atendido exatamente uma vez;
  \item cada rota respeita capacidade $Q$;
  \item cada atendimento satisfaz $t^{start}_i \le b_i + \varepsilon$;
  \item o número de rotas não ultrapassa $m$.
\end{enumerate}

O objetivo adotado é minimizar custo total de deslocamento:
\[
 \min \sum_{r\in R}\sum_{(i,j)\in r} d_{ij}.
\]

\subsection{Solomon I1 implementado}
\label{subsec:solomon_i1_impl}

A solução inicial usada em todo o pipeline é produzida por uma versão clássica do I1 com seed determinística e inserção paralela sobre as rotas abertas.
As escolhas fixas do trabalho são:
\[
\alpha_1=1,\quad \alpha_2=0,\quad \lambda=1,\quad \mu=1.
\]

Mesmo com $\alpha_2=0$, a verificação temporal continua obrigatória: o termo temporal deixa de pesar no escore, mas continua atuando como restrição de factibilidade.

\subsubsection{Critério de inserção $c_1$}

Ao testar inserção do cliente $u$ entre nós adjacentes $(i,j)$ de uma rota, define-se:
\begin{align}
 c_{11}(i,u,j) &= d_{iu} + d_{uj} - \mu d_{ij}, \\
 c_{12}(i,u,j) &= b_j^{new} - b_j^{old}, \\
 c_1(i,u,j) &= \alpha_1 c_{11}(i,u,j) + \alpha_2 c_{12}(i,u,j).
\end{align}

Na implementação, $b_j^{old}$ e $b_j^{new}$ são os tempos de início de serviço de $j$ antes e depois da inserção, obtidos pelo recálculo do cronograma da rota base e da rota candidata.

\subsubsection{Critério de seleção global $c_2$}

Para cada cliente não roteirizado $u$, primeiro identifica-se sua melhor inserção (menor $c_1$).
Depois calcula-se:
\[
 c_2(u)=\lambda d_{0u}-c_1(u).
\]
O cliente inserido na iteração é o de maior $c_2$.
Em empate de $c_2$, a implementação desempata por menor $c_1$.

\subsubsection{Política de abertura de rota}

Se nenhuma inserção factível for encontrada nas rotas atuais, o algoritmo abre nova rota com o cliente mais distante do depósito (\emph{farthest seed}), desde que ainda exista veículo disponível.

\begin{algorithm}[htbp]
\caption{Solomon I1}
\label{alg:i1}
\begin{algorithmic}[1]
\Require instância VRPTW e parâmetros $(\alpha_1,\alpha_2,\lambda,\mu,\varepsilon)$
\Ensure solução inicial $S$
\State $S\leftarrow\emptyset$, $NR\leftarrow N$
\While{$NR\neq\emptyset$}
  \If{$S=\emptyset$}
    \State seleciona seed mais distante do depósito e abre rota $(0,u,0)$
    \State remove $u$ de $NR$
    \State \textbf{continue}
  \EndIf
  \State busca melhor inserção global (todas as rotas, posições e clientes em $NR$)
  \If{existe inserção factível}
    \State escolhe $u^*=\argmax c_2(u)$ e aplica inserção correspondente
    \State $NR\leftarrow NR\setminus\{u^*\}$
  \Else
    \If{$|S|<m$}
      \State abre nova rota com seed farthest
      \State remove seed de $NR$
    \Else
      \State \textbf{break}
    \EndIf
  \EndIf
\EndWhile
\State recomputa $S$ e remove rotas triviais
\State \Return $S$
\end{algorithmic}
\end{algorithm}

\paragraph{Texto complementar ao pseudocódigo.}
A etapa da linha 8 (``melhor inserção global'') é o núcleo da implementação: para cada cliente não roteirizado, o código percorre todas as rotas e todas as posições de inserção, computa o menor $c_1$ factível desse cliente e só então aplica o desempate global por maior $c_2$.
Isso evita decisões miopes por rota e aproxima o comportamento descrito por Solomon para o I1 clássico.

No nível de estrutura de dados, cada rota é representada por uma sequência ordenada de clientes com custo, carga e cronograma associados.
Quando um candidato $(i,u,j)$ é testado, o algoritmo cria uma rota temporária, insere $u$, executa \texttt{recompute} e valida simultaneamente capacidade e janelas.
Somente candidatos factíveis entram na competição por $c_1$ e $c_2$.
Se nenhum candidato for viável, a linha 14 abre nova rota (se houver veículo disponível); caso contrário, a construção encerra.

\subsection{Papel do I1 no estudo comparativo}

Neste trabalho, o I1 é o baseline comum para VND, GRASP e Tabu.
Essa decisão metodológica separa claramente duas camadas:
\begin{itemize}
  \item qualidade da construção inicial;
  \item capacidade de melhoria da busca local/meta-heurística.
\end{itemize}

Com isso, diferenças de desempenho observadas entre algoritmos são atribuídas principalmente ao mecanismo de busca, e não a inicializações heterogêneas.
